\documentclass[final]{beamer}
\usepackage[orientation=portrait,size=a0,scale=1.15]{beamerposter}
\usepackage[utf8]{inputenc}
\usepackage{amsmath,amssymb}
\usepackage{tikz}
\usetikzlibrary{shapes.geometric,arrows.meta,positioning,calc,fit}
\usepackage{xcolor}
\usepackage{booktabs}
\usepackage{multicol}
\usepackage{tcolorbox}
\tcbuselibrary{skins}
\usepackage{array}
\usepackage{colortbl}
\usepackage{pgfplots}
\pgfplotsset{compat=1.18}

% ── Colors ───────────────────────────────────────────────────────────────────
\definecolor{primary}{HTML}{293AA2}
\definecolor{secondary}{HTML}{FFCD00}
\definecolor{darkprimary}{HTML}{1A2575}
\definecolor{lightprimary}{HTML}{DDE1F5}
\definecolor{darkgray}{HTML}{333333}
\definecolor{accentgray}{HTML}{777777}
\definecolor{rowA}{HTML}{E8EBFA}
\definecolor{rowB}{HTML}{FFFFFF}

% ── Beamer setup ──────────────────────────────────────────────────────────────
\usetheme{default}
\setbeamercolor{structure}{fg=primary}
\setbeamercolor{block title}{bg=primary,fg=white}
\setbeamercolor{block body}{bg=lightprimary,fg=darkgray}
\setbeamerfont{block title}{size=\large,series=\bfseries}
\setbeamerfont{block body}{size=\normalsize}
\setbeamertemplate{navigation symbols}{}

% ── tcolorbox styles ──────────────────────────────────────────────────────────
\tcbset{
  pbox/.style={enhanced,arc=5pt,colback=lightprimary,colframe=primary,
    boxrule=2pt,left=8pt,right=8pt,top=6pt,bottom=6pt,drop shadow={opacity=0.18}},
  ybox/.style={enhanced,arc=5pt,colback=secondary!28,colframe=secondary!65!darkgray,
    boxrule=1.8pt,left=7pt,right=7pt,top=5pt,bottom=5pt},
  dbox/.style={enhanced,arc=5pt,colback=primary,colframe=darkprimary,
    boxrule=0pt,left=9pt,right=9pt,top=7pt,bottom=7pt,drop shadow},
  pboxtight/.style={enhanced,arc=4pt,colback=lightprimary,colframe=primary,
    boxrule=1.5pt,left=6pt,right=6pt,top=4pt,bottom=4pt}
}

% ── Table helpers ─────────────────────────────────────────────────────────────
\newcommand{\thead}[1]{\cellcolor{primary}\color{white}\textbf{#1}}
\newcommand{\trA}[1]{\cellcolor{rowA}#1}
\newcommand{\trB}[1]{\cellcolor{rowB}#1}

% ── Header ───────────────────────────────────────────────────────────────────
\setbeamertemplate{headline}{%
\begin{beamercolorbox}[wd=\paperwidth,ht=13cm,dp=0.5cm]{block title}
\begin{tikzpicture}[remember picture,overlay]
  \fill[darkprimary](0,0)rectangle(\paperwidth,14cm);
  \fill[primary](0,0.7cm)rectangle(\paperwidth,13.3cm);
  \fill[secondary](0,0)rectangle(\paperwidth,0.7cm);
  \fill[secondary](0,13.3cm)rectangle(\paperwidth,14cm);
  %% ── Wide banner scene: car drives left→right along a winding track ──────────
  %% All coords in cm, origin = bottom-left of header box.
  %% The path winds across the full ~84 cm width at ~7 cm height (mid-banner).

  %% --- Centerline path (the car follows this) ---
  %% Starts at x≈3, winds: slight dip, rise, dip, rise to exit right edge
\draw[white,opacity=0.22,line width=5pt] 
  plot[smooth,tension=0.55] coordinates {
    (2,6.4) (6,6.35) (10,5.8) (14,5.9)  (22,7.5) 
    (30,5.35)  (38,5.4) (46,7.9)
    (54,6.3) (62,5.35) (70,6.4) (74,7.1)
     (82,7.5)
  };


  %% --- Macro for a cone triangle (top-view, side-lit) ---
  %% Blue cones along upper boundary (above centerline ~+2.5 cm)
  %% Yellow cones along lower boundary (below centerline ~-2.5 cm)
  %% We place them by hand along the winding path

  %% BLUE cones (upper row)
  \foreach \cx/\cy in {
    2/9.8, 6/9.5, 10/8.8, 14/8.3, 18/10.2, 22/10.6, 26/10.4,
    30/8.5, 34/8.0, 38/8.0, 42/10.8, 46/11.0, 50/10.7,
    54/8.9, 58/8.0, 62/8.2, 66/8.5, 70/9.3, 74/10.2, 78/10.5, 82/10.8}{
    \begin{scope}[shift={(\cx,\cy)}]
        % square base
        \fill[primary!80] (-0.35,0) rectangle (0.35,0.12);

        % cone body
        \fill[primary!70] (-0.22,0.12) -- (0.22,0.12) -- (0,1.1) -- cycle;

        % white reflective strip
        \fill[white] (-0.12,0.55) -- (0.12,0.55) -- (0.07,0.75) -- (-0.07,0.75) -- cycle;
    \end{scope}
  }

  %% YELLOW cones (lower row)
  \foreach \cx/\cy in {
    4/3.0, 8/3.2, 12/2.8, 16/3.5, 20/4.2, 24/4.4,
    28/2.6, 32/2.2, 36/2.0, 40/2.8, 44/4.5, 48/4.8,
    52/4.2, 56/2.8, 60/2.3, 64/2.5, 68/2.8, 72/3.5, 76/4.0, 80/4.2}{
    \begin{scope}[shift={(\cx,\cy)}]
        % square base
        \fill[secondary!80] (-0.35,0) rectangle (0.35,0.12);

        % cone body
        \fill[secondary!70] (-0.22,0.12) -- (0.22,0.12) -- (0,1.1) -- cycle;

        % white reflective strip
        \fill[black] (-0.12,0.55) -- (0.12,0.55) -- (0.07,0.75) -- (-0.07,0.75) -- cycle;
    \end{scope}
  }

  %% --- Car: exact user-specified graphic, placed at left of banner ---
  \begin{scope}[shift={(6.5cm,6cm)},scale=1.2]
    \fill[secondary!85](-2,-0.8)--(-1.5,0.9)--(1.5,0.9)--(2,-0.8)--cycle;
    \fill[white!15!primary](-1.1,0.9)--(-0.7,2.0)--(0.7,2.0)--(1.1,0.9)--cycle;
    \fill[white](-1.5,-1.0)circle(0.33); \fill[white](1.5,-1.0)circle(0.33);
    \fill[darkprimary](-1.5,-1.0)circle(0.16); \fill[darkprimary](1.5,-1.0)circle(0.16);
    \foreach \a in {15.00, 11.00, 8.00, 5.00, 3.00, 2.00, 1.67, 1.33, 1.00, 0.67, 0.33, 0.00, -0.33, -0.67, -1.00, -1.33, -1.67, -2.00, -2.33, -2.67, -3.00, -3.33, -3.67, -4.00, -4.33, -4.67, -5.00, -5.33, -5.67, -6.00, -7.00, -8.00, -9.00, -10.00, -11.00, -12.00, -13.00, -14.00, -19.00, -25.00}
      {\draw[secondary,opacity=0.5,line width=0.6pt](0,1.4)--+(\a:9.9);}
    \fill[secondary](0,1.4)circle(0.12);
  \end{scope}
  %% Title text
  \node[white,font=\fontsize{60pt}{68pt}\selectfont\bfseries,anchor=center]
    at(0.5\paperwidth,10.8cm){Lund Formula Student    Driverless};
  \node[secondary,font=\fontsize{31pt}{38pt}\selectfont,anchor=center]
    at(0.5\paperwidth,8.3cm){LiDAR-Based Autonomous Racing: Perception $\cdot$ SLAM $\cdot$ Path Planning $\cdot$ Tracking};
\end{tikzpicture}
\end{beamercolorbox}}
% ── Footer ────────────────────────────────────────────────────────────────────
\setbeamertemplate{footline}{%
\begin{beamercolorbox}[wd=1.0\paperwidth,ht=3.0cm,dp=0.4cm]{block title}
\begin{tikzpicture}[remember picture,overlay]
  \fill[secondary](0,0)rectangle(\paperwidth,3.4cm);
  \node[darkprimary,font=\normalsize\bfseries]at(0.5\paperwidth,2.1cm){%
    Alve Lindell (Manager) $\cdot$
    William Andersson (SW) $\cdot$
    Felix Hellborg (SW) $\cdot$
    Ahmad Hamid (SW) $\cdot$
    Lucas Bertman (SW) $\cdot$
    Tommaso Batori (Controls) $\cdot$
    Diego Loayza (Controls)};
\end{tikzpicture}
\end{beamercolorbox}}

\begin{document}
\begin{frame}[t,fragile]{}
\vspace{0.7cm}

%% ── Overview banner ─────────────────────────────────────────────────────────
\begin{tcolorbox}[dbox,width=\textwidth]
\begin{columns}[c]
\begin{column}{0.48\textwidth}
  \color{white}\large\bfseries Driverless Events \& Actuation\normalfont\\[5pt]
  \color{white}\normalsize
  Same dynamic events as driven --- without a driver.
  \textbf{Trackdrive} replaces Endurance: 10 laps fully autonomous.
  The car controls steering, braking, and throttle via electronic
  actuation, using only \textbf{LiDAR} as its autonomous sensor.
\end{column}
\begin{column}{0.22\textwidth}
  \centering
  \begin{tcolorbox}[ybox,left=5pt,right=5pt,top=4pt,bottom=4pt]
    \color{darkprimary}\centering\bfseries Autonomous Actuation\\[4pt]
    \normalfont\small
    Steering\\Braking\\Throttle (electronic)
  \end{tcolorbox}
\end{column}
\begin{column}{0.28\textwidth}
  \centering
  \begin{tikzpicture}[node distance=0.5cm,scale=0.85]
    \tikzstyle{pn}=[fill=secondary,draw=darkprimary,rounded corners=3pt,
      text width=2.6cm,align=center,minimum height=0.58cm,
      font=\small\bfseries,text=darkprimary]
    \tikzstyle{bn}=[fill=white!18!primary,draw=white!40,rounded corners=3pt,
      text width=2.6cm,align=center,minimum height=0.58cm,
      font=\small\bfseries,text=white]
    \node[pn](s){LiDAR};
    \node[bn,right=0.3cm of s](p){Perception};
    \node[bn,right=0.3cm of p](sl){Mapping};
    \node[bn,right=0.3cm of sl](pl){Path Plan};
    \node[pn,right=0.3cm of pl](m){Tracking};
    \foreach \a/\b in {s/p,p/sl,sl/pl,pl/m}
      {\draw[-{Stealth[length=6pt]},white,line width=1.5pt](\a)--(\b);}
  \end{tikzpicture}
\end{column}
\end{columns}
\end{tcolorbox}

\vspace{0.6cm}

%% ── Row 1: Hardware | LiDAR Sensor | LIO/SLAM ───────────────────────────────
\begin{columns}[t,totalwidth=\textwidth]

%% ─ LIDAR SENSOR ──────────────────────────────────────────────────────────────
\begin{column}{0.248\textwidth}

\begin{beamercolorbox}[rounded=true,sep=0.30cm,wd=\linewidth]{block title}
  \large\bfseries Main Sensors

\end{beamercolorbox}

\begin{beamercolorbox}[rounded=true,sep=0.38cm,wd=\linewidth]{block body}
\centering

  \begin{tcolorbox}[pbox]
    {\color{primary}\bfseries Hesai Pandar40P}\\[3pt]
    \small
    $\bullet$ 40 channels, 360° horizontal FoV, 10 Hz scan rate\\
    $\bullet$ Range up to 200 m, with high accuracy and resolution\\
    $\bullet$ Vertical FoV concentrated around -6° to +2° degrees \to perfect for seeing cones\\
    $\bullet$ Horizontal FoV partially obstructed by nose cone $\to$ only ~170° usable\\
    $\bullet$ No built-in IMU $\to$ separate INS unit\\
    $\bullet$ GPS timestamp synchronization for accurate sensor fusion\\
    $\bullet$ Mounted at the front of the car \to maximum visability
  \end{tcolorbox}
  \begin{tikzpicture}[scale=0.82]
    \draw[primary,opacity=0.10,line width=1pt](0,0)ellipse(4.0cm and 1.52cm);
    \draw[primary,opacity=0.18,line width=1pt](0,0)ellipse(3.3cm and 1.25cm);
    \draw[primary,opacity=0.30,line width=1pt](0,0)ellipse(2.6cm and 0.99cm);
    \draw[primary,opacity=0.45,line width=1pt](0,0)ellipse(2.0cm and 0.76cm);
    \draw[primary,opacity=0.62,line width=1pt](0,0)ellipse(1.4cm and 0.53cm);
    \draw[primary,opacity=0.80,line width=1pt](0,0)ellipse(0.8cm and 0.30cm);
    \fill[darkprimary](-0.34,-0.34)rectangle(0.34,0.34);
    \fill[secondary](-0.20,0.34)rectangle(0.20,0.82);
    % Only front 90 deg (nose cone constraint)
    \foreach \a in {-85,-70,-55,-40,-25,-10,0,10,25,40,55,70,85}
        {\draw[secondary,opacity=0.65,line width=1.2pt](0,0.1)--+(\a:3.6);}
    \foreach \a in {-180,-165,...,-95,95,110,...,180}
        {\draw[primary!40,opacity=0.25,line width=0.4pt](0,0.1)--+(\a:3.6);}
    % Shade used FOV
    \fill[secondary,opacity=0.08]
      (0,0.1)--(-45:3.6)arc(-45:45:3.6)--(0,0.1);
    \node[secondary!70!darkgray,font=\scriptsize\bfseries]at(0,2.8){Used FOV (~170\textdegree)};
    \node[accentgray,font=\scriptsize]at(-3.2,0.6){Unused};
    \node[accentgray,font=\small]at(0,-1.65){Top view --- 40-channel scan rings};
  \end{tikzpicture}

  \begin{tcolorbox}[pbox]
    {\color{primary}\bfseries Movella Xsense MTI-680G}\\[3pt]
    \small
    $\bullet$ Buliltin 9-axis IMU\\
    $\bullet$ High-rate (up to 400 Hz) inertial measurements\\
    $\bullet$ GPS antenna for timestamp synchronization\\
    $\bullet$ RTK for ground truth validation
  \end{tcolorbox}
  \begin{tikzpicture}[>=Stealth, font=\sffamily\small, scale=1.9]

    \definecolor{imucolor}{RGB}{255,248,210}
    \definecolor{axisblue}{RGB}{30,100,200}
    \definecolor{axisred}{RGB}{200,40,40}
    \definecolor{axisgreen}{RGB}{20,150,60}
    \definecolor{magpurple}{RGB}{140,40,180}
    \definecolor{gpsgreen}{RGB}{30,150,80}

    %% IMU box
    \draw[thick, fill=imucolor, rounded corners=4pt]
        (-0.7,-0.6) rectangle (0.7,0.6);

    %% Body frame axes
    \draw[->, very thick, axisblue]
        (0,0) -- (3.0,0) node[right] {$x$};
    \draw[->, very thick, axisgreen]
        (0,0) -- (0,2.6) node[above] {$z$};
    \draw[->, very thick, axisred]
        (0,0) -- (-0.9,-1.5) node[below left] {$y$};

    %% Linear accelerations (dashed)
    \draw[->, thick, axisblue, dashed]
        (0,0) -- (1.5,0) node[below] {$a_x$};
    \draw[->, thick, axisgreen, dashed]
        (0,0) -- (0,1.4) node[right] {$a_z$};
    \draw[->, thick, axisred, dashed]
        (0,0) -- (-0.45,-0.75) node[left] {$a_y$};

    %% --- Angular velocities as circular arrows ---
    \draw[axisred!80, thick, ->]
        (-0.4,-1.4)
        arc[start angle=0,end angle=270,
            x radius=0.45, y radius=0.45];

    \node[axisred!80, font=\scriptsize\bfseries]
        at (-1.8,-0.8) {Pitch};

    \draw[axisblue!80, thick]
        (2.6,0.55) arc[start angle=90,end angle=270,
        x radius=0.4, y radius=0.55];
    \draw[axisblue!80, thick, ->]
        (2.6,-0.55) arc[start angle=-90,end angle=-20,
        x radius=0.4, y radius=0.55];
    \node[axisblue!80, font=\scriptsize\bfseries]
        at (3.15,0.8) {Roll};

    \draw[axisgreen!80, thick]
        (-0.55,2.2) arc[start angle=180,end angle=360,
        x radius=0.55, y radius=0.4];
    \draw[axisgreen!80, thick, ->]
        (0.55,2.2) arc[start angle=0,end angle=40,
        x radius=0.55, y radius=0.4];
    \node[axisgreen!80, font=\scriptsize\bfseries]
        at (-1.2,2.4) {Yaw};

    %% --- Magnetometer vector (cleaner placement) ---
    \draw[->, thick, magpurple, dotted]
        (0,0) -- (-1.2,0.9);
    \node[magpurple, font=\scriptsize] at (-1.3,1.0)
        {$\vec{B}$};

    %% --- GPS antenna ---
    \draw[thick, black!60] (1.8, 0.4) -- (1.8, 1.5);
    \draw[thick, fill=gpsgreen!30] (1.5,1.5) -- (2.1,1.5) --
        (2.0,1.85) -- (1.6,1.85) -- cycle;

    \foreach \r in {0.3, 0.55, 0.80}{
        \draw[orange!80]
            ($(1.8,1.95)+(-\r*0.5, \r*0.3)$)
            arc (145:35:\r*0.6);
    }
    \node[gpsgreen!70!black, font=\scriptsize] at (1.8,0.2) {GPS};

    %% Rate label
    \node[fill=yellow!15, draw=black!30, rounded corners=3pt,
        font=\scriptsize, inner sep=3pt] at (0,-2.2)
        {\textbf{400 Hz} inertial measurements};

  \end{tikzpicture}
  \begin{tcolorbox}[pbox]
    {\color{primary}\bfseries Other Sensors}\\[3pt]
    \small
    $\bullet$ 4x Wheel speed\\
    $\bullet$ Steering angle sensor\\
    $\bullet$ Torque at each wheel\\
    $\bullet$ Brake pressure sensor
  \end{tcolorbox}
  \begin{tcolorbox}[ybox]
    \small\color{darkprimary}
    \textbf{Rationale and Caveats:} Using LiDAR only gives very accurate cone detection, localization and
    allows us to run the entier pipeline on CPU. However, it means that reliable color detection is not possible,
    which makes path planning signicantly more challenging. Adding a camera would allow us to also detect cone colors,
    but due to the accuracy of our pathplanning the complexity added by a camera would not be worth it.
  \end{tcolorbox}
  \end{beamercolorbox}
\end{column}
%% ─ PERCEPTION ────────────────────────────────────────────────────────────────
\begin{column}{0.248\textwidth}
\begin{beamercolorbox}[rounded=true,sep=0.30cm,wd=\linewidth]{block title}
  \large\bfseries Perception Pipeline
\end{beamercolorbox}
\begin{beamercolorbox}[rounded=true,sep=0.38cm,wd=\linewidth]{block body}
  \centering

  \begin{tcolorbox}[pbox]
    {\color{primary}\bfseries Ground Removal --- LineFit}\\[3pt]
    \small
    Bins point cloud in \textbf{polar coordinates},
    then uses \textbf{RANSAC} to fit lines between bins.
    Ground points are removed; remainder is the obstacle cloud.\\
    $\bullet$ Compared to simpler ground removal this can handle non planar ground
  \end{tcolorbox}
\vspace{0.9cm}
\centering
\begin{tikzpicture}[scale=2.58]

    % Ground line
    \draw[secondary,line width=2pt](-3.5,0.3)--(3.5,0.3);

    % Ground points
    \foreach \x in {-3.5,-3.4,...,3.5} {
        \pgfmathsetmacro{\dx}{(rand*0.04-0.02)}
        \pgfmathsetmacro{\dy}{(rand*0.02-0.01)}
        \fill[accentgray,opacity=0.5] (\x+\dx,0.3+\dy) circle(0.05);
    }

    % Function to draw a LiDAR-style cone with noise
    % #1: x-center, #2: base y, #3: height, #4: base half-width, #5: number of scan lines, #6: color
    \newcommand{\lidarcone}[6]{%
    \foreach \i in {0,...,#5} {%
        \pgfmathsetmacro{\y}{#2 + \i*(#3/#5)} % height of this scan line
        \pgfmathsetmacro{\halfwidth}{#4*(#5-\i)/#5} % taper width
        % number of points proportional to width
        \pgfmathtruncatemacro{\npts}{int(2*\halfwidth*10)+1}
        \foreach \j in {0,...,\npts} {%
        \pgfmathsetmacro{\xp}{#1 - \halfwidth + \j*2*\halfwidth/\npts}
        % add random noise
        \pgfmathsetmacro{\dx}{(rand*0.02-0.01)}
        \pgfmathsetmacro{\dy}{(rand*0.02-0.01)}
        \fill[#6] (\xp+\dx,\y+\dy) circle(0.07);
        }%
    }%
    }

    % Draw cones with noise
    \lidarcone{-2.3}{0.5}{0.5}{0.2}{2}{primary}   % primary cone 1
    \lidarcone{0.6}{0.5}{0.5}{0.2}{2}{primary}    % primary cone 2
    \lidarcone{1.8}{0.5}{0.5}{0.2}{2}{secondary}  % secondary cone

    % Labels
    \node[accentgray,font=\scriptsize] at (0,0.15) {Ground (removed)};

\end{tikzpicture}

  \vspace{0.3cm}
  \begin{tcolorbox}[pbox]
    {\color{primary}\bfseries Clustering}\\[3pt]
    \small
    Clusters obstacle cloud to extract distinct objects after removing the ground.\\[3pt]
    Algorithm: \textbf{Euclidean clustering}\\
    $\bullet$ Cluster by Euclidean distance\\
    $\bullet$ Efficient for large point clouds\\
    $\bullet$ Works well for well-separated objects like cones\\
    $\bullet$ DBSCAN is unnecessary due to clear cone shapes and worse performance\\
    \textbf{Cone classification} by shape constraints:\\
    $\bullet$ Height: 2--40 cm\\
    $\bullet$ Radius: 17 cm\\
    $\bullet$ Any cluster that fits these constraints is classified as a cone\\
    $\bullet$ False positives are filtered out in SLAM since multiple observations of the same cone are required to create a landmark
  \end{tcolorbox}

  \vspace{0.3cm}
  % Clustering diagram
  \begin{tikzpicture}[scale=2.9]
    % Cluster 1 (blue)
    \foreach \x/\y in {-2.4/0.8,-2.2/1.1,-2.0/0.9,-2.3/1.3,-2.1/0.7}
      {\fill[primary,opacity=0.85](\x,\y)circle(0.09);}
    \draw[primary,line width=1.2pt,dashed,rounded corners=8pt]
      (-2.6,0.55)rectangle(-1.8,1.5);
    \node[primary,font=\scriptsize\bfseries]at(-2.2,1.7){Cluster A};
    % Cluster 2 (yellow)
    \foreach \x/\y in {1.5/0.9,1.7/1.2,1.9/0.8,1.6/1.4,2.0/1.1}
      {\fill[secondary,opacity=0.85](\x,\y)circle(0.09);}
    \draw[secondary!70!darkgray,line width=1.2pt,dashed,rounded corners=8pt]
      (1.3,0.55)rectangle(2.2,1.6);
    \node[secondary!70!darkgray,font=\scriptsize\bfseries]at(1.75,1.8){Cluster B};
    % Noise
    \foreach \x/\y in {-0.5/1.5,0.3/0.7,-1.0/0.5,0.8/1.8}
      {\fill[accentgray,opacity=0.5](\x,\y)circle(0.07);}
    \node[accentgray,font=\scriptsize]at(-0.1,0.3){noise};
  \end{tikzpicture}

  \vspace{0.3cm}
  \begin{tcolorbox}[ybox]
    \small\color{darkprimary}
    \textbf{Challenge:} Noisy measurements, partial observations, noisy pose
    $\to$ need \textbf{SLAM} to find ``real'' cone positions
  \end{tcolorbox}

\end{beamercolorbox}
\end{column}

%% ─ LIO / SLAM ────────────────────────────────────────────────────────────────
\begin{column}{0.248\textwidth}
\begin{beamercolorbox}[rounded=true,sep=0.30cm,wd=\linewidth]{block title}
  \large\bfseries LIO \& GraphSLAM
\end{beamercolorbox}
\begin{beamercolorbox}[rounded=true,sep=0.38cm,wd=\linewidth]{block body}

  \begin{tcolorbox}[pbox]
    {\color{primary}\bfseries Direct LiDAR-Inertial Odometry (DLIO)}\\[3pt]
    \small
    LiDAR-Inertial Odometry accomplishes two things simultaneously:
    SLAM on the point cloud to accumulate scans, and accurate
    high-rate odometry for vehicle state estimation.\\[4pt]
    \textbf{Key steps:}\\
    $\bullet$ Continuous IMU integration + \textbf{deskewing}\\
    $\bullet$ \textbf{GICP} scan-to-map matching (IMU prior)\\
    $\bullet$ Geometric observer for sensor fusion
  \end{tcolorbox}

  \vspace{0.3cm}
  \begin{tcolorbox}[ybox]
    {\color{darkprimary}\bfseries DLIO vs FAST-LIO}\\[3pt]
    \small
    \textbf{Pros:} More stable; significantly lower CPU cost\\
    \textbf{Cons:} Lag-sensitive --- if LiDAR misses a publish, pipeline breaks
  \end{tcolorbox}

  \vspace{0.3cm}
  % DLIO pipeline
  \centering
  \begin{tikzpicture}[node distance=0.48cm,scale=0.9]
    \tikzstyle{dn}=[fill=primary,text=white,rounded corners=3pt,
      minimum width=4.0cm,minimum height=0.55cm,align=center,font=\small\bfseries]
    \tikzstyle{yn}=[fill=secondary,text=darkprimary,rounded corners=3pt,
      minimum width=4.0cm,minimum height=0.55cm,align=center,font=\small\bfseries]
    \node[yn](imu){IMU @ high rate};
    \node[dn,below=0.38cm of imu](desk){Deskewing};
    \node[dn,below=0.38cm of desk](gicp){GICP matching};
    \node[yn,below=0.38cm of gicp](map){Local Map Update};
    \node[dn,below=0.38cm of map](odom){Odometry Output};
    \foreach \a/\b in {imu/desk,desk/gicp,gicp/map,map/odom}
      {\draw[-{Stealth[length=5pt]},primary,line width=1.2pt](\a)--(\b);}
    \draw[secondary!70!darkgray,dashed,line width=0.8pt]
      (imu.east)--++(0.55,0)|-(gicp.east);
    \node[accentgray,font=\scriptsize,rotate=90]at(2.75,-2.1){IMU prior};
  \end{tikzpicture}

  \vspace{0.35cm}
  \begin{tcolorbox}[pbox]
    {\color{primary}\bfseries GraphSLAM with ISAM2 (GTSAM)}\\[3pt]
    \small
    $\bullet$ Solves SLAM \textbf{globally} $\to$ global minima, less drift\\
    $\bullet$ \textbf{Loop closure} corrects accumulated error\\
    $\bullet$ ISAM2 exploits sparsity + incremental updates\\
    \quad for real-time performance
  \end{tcolorbox}

  \vspace{0.3cm}
  % SLAM map mini diagram
  \begin{tikzpicture}[scale=0.85]
    \draw[primary!25,line width=7pt,rounded corners=18pt]
      (-2.5,0)--(-1.7,1.3)--(-0.3,2.1)--(1.3,2.1)--(2.2,1.3)--(2.5,0)
      --(2.2,-1.3)--(1.3,-2.1)--(-0.3,-2.1)--(-1.7,-1.3)--cycle;
    \draw[secondary!45,line width=3.5pt,rounded corners=11pt]
      (-1.6,0)--(-1.0,0.9)--(-0.1,1.4)--(1.0,1.4)--(1.6,0.9)--(1.9,0)
      --(1.6,-0.9)--(1.0,-1.4)--(-0.1,-1.4)--(-1.0,-0.9)--cycle;
    \foreach \x/\y in {-2.5/0,-1.7/1.3,-0.3/2.1,1.3/2.1,2.2/1.3,2.5/0,
                        2.2/-1.3,1.3/-2.1,-0.3/-2.1,-1.7/-1.3}
      {\fill[primary](\x,\y)circle(0.12);}
    \foreach \x/\y in {-1.6/0,-1.0/0.9,-0.1/1.4,1.0/1.4,1.6/0.9,1.9/0,
                        1.6/-0.9,1.0/-1.4,-0.1/-1.4,-1.0/-0.9}
      {\fill[secondary](\x,\y)circle(0.11);}
    \fill[darkprimary](0,0)circle(0.17);
    \draw[-{Stealth},red!50!black,line width=1.1pt,dashed]
      (-0.7,0.2)arc[start angle=140,end angle=40,radius=0.9];
    \node[red!50!black,font=\scriptsize]at(0.2,0.65){LC};
    \node[accentgray,font=\small]at(0,-2.6){Cone landmark map};
  \end{tikzpicture}

\end{beamercolorbox}
\end{column}


\end{columns}

\vspace{0.6cm}

%% ── Row 2: Path Planning | Tracking | Simulator ─────────────────────────────
\begin{columns}[t,totalwidth=\textwidth]

%% ─ PATH PLANNING ─────────────────────────────────────────────────────────────
\begin{column}{0.48\textwidth}

\begin{beamercolorbox}[rounded=true,sep=0.30cm,wd=\linewidth]{block title}
  \large\bfseries Path Planning
\end{beamercolorbox}

\begin{beamercolorbox}[rounded=true,sep=0.38cm,wd=\linewidth,ht=37.3cm]{block body}

\begin{columns}[t]
\begin{column}{0.48\linewidth}

  \begin{tcolorbox}[pbox]
    {\color{primary}\bfseries Trackdrive: Graph Search}\\[3pt]
    \small
    Idea: Centerline $\in$ Midpoints crossing track:\\
    \textbf{Delaunay triangulate} all cone positions\\
    $\bullet$ Filter on edge length, angles, adjacency\\
    $\bullet$ Build search graph from adjacent midpoints \\
    $\bullet$ Weight edges by quality (correct lengths, angle, etc.)\\
    \textbf{Dijkstra} $\to$ path cost \\
    Longest path with lowest cost $\to$ centerline
  \end{tcolorbox}

  \vspace{0.9cm}
  \centering
  \input{track.tex}

\end{column}

\begin{column}{0.48\linewidth}

  \begin{tcolorbox}[pbox]
    {\color{primary}\bfseries Acceleration: RANSAC Line Fit}\\[3pt]
    \small
    Fit two lines (one per cone row) using \textbf{RANSAC}:\\
    $\bullet$ Robust to outlier cones\\
    $\bullet$ Utilizes the simplicity of the problem\\
    Centerline = average of both lines
  \end{tcolorbox}

  \vspace{0.3cm}
  \begin{tcolorbox}[pbox]
    {\color{primary}\bfseries Skidpad: Optimization}\\[3pt]
    \small
    Known path $\to$ formulate as \textbf{optimization}:\\
    $\bullet$ \textbf{Tukey loss} for outlier robustness\\
    $\bullet$ Minimize residuals vs. circular arcs
  \end{tcolorbox}

  \vspace{8.8cm}
  % Skidpad circular path
  \centering

  \begin{tikzpicture}[scale=0.46]

    % Skidpad parameters
    \def\innerDia{15.25}    % inner circle diameter
    \def\outerDia{21.25}    % outer circle diameter
    \def\radiusInner{7.625} % radius
    \def\radiusOuter{10.625}
    \def\centerSep{18.25}   % distance between circle centers

    % Centers
    \coordinate (Cleft)  at (0,0);
    \coordinate (Cright) at (\centerSep,0);

    % --- Function to add noise ---
    \newcommand{\noisycone}[2]{% #1 = center, #2 = color
        % Box-Muller: generate gaussian with mean 0, std dev 0.2
        \pgfmathsetmacro{\u}{rand}%
        \pgfmathsetmacro{\v}{rand}%
        \pgfmathsetmacro{\r}{sqrt(-2*ln(\u))*0.2}%
        \pgfmathsetmacro{\theta}{2*3.14159265*\v}%
        \pgfmathsetmacro{\dx}{\r*cos(\theta)}%
        \pgfmathsetmacro{\dy}{\r*sin(\theta)}%
        \fill[#2] ($ (#1) + (\dx,\dy) $) circle (0.3);
    }


    % --- Draw noisy cones ---
    \foreach \i in {0,...,16} { % 17 inner cones
        \pgfmathsetmacro{\ang}{\i*360/17}
        \pgfmathsetmacro{\rad}{\radiusInner}
        % left inner
        \pgfmathsetmacro{\noiseX}{(\rad + (rand*0.5-0.25))*cos(\ang)}
        \pgfmathsetmacro{\noiseY}{(\rad + (rand*0.5-0.25))*sin(\ang)}
        \fill[primary] ($(Cleft) + (\noiseX,\noiseY)$) circle (0.4);

        % right inner
        \pgfmathsetmacro{\noiseX}{(\rad + (rand*0.5-0.25))*cos(\ang)}
        \pgfmathsetmacro{\noiseY}{(\rad + (rand*0.5-0.25))*sin(\ang)}
        \fill[secondary] ($(Cright) + (\noiseX,\noiseY)$) circle (0.4);
    }

    \foreach \i in {2,...,11} { % 13 outer cones
        \pgfmathsetmacro{\ang}{\i*360/13}
        \pgfmathsetmacro{\rad}{\radiusOuter}
        % left outer
        \pgfmathsetmacro{\noiseX}{(\rad + (rand*0.5-0.25))*cos(\ang)}
        \pgfmathsetmacro{\noiseY}{(\rad + (rand*0.5-0.25))*sin(\ang)}
        \fill[secondary] ($(Cleft) + (\noiseX,\noiseY)$) circle (0.4);

    }
    \foreach \i in {2,...,11} { % 13 outer cones
        \pgfmathsetmacro{\ang}{180+\i*360/13}
        \pgfmathsetmacro{\rad}{\radiusOuter}

        % right outer
        \pgfmathsetmacro{\noiseX}{(\rad + (rand*0.5-0.25))*cos(\ang)}
        \pgfmathsetmacro{\noiseY}{(\rad + (rand*0.5-0.25))*sin(\ang)}
        \fill[primary] ($(Cright) + (\noiseX,\noiseY)$) circle (0.4);
    }

    % --- Draw best-fit figure-eight arcs ---
    \draw[ultra thick, red!70] 
        (Cleft) circle (\radiusInner); % left inner
    \draw[ultra thick, red!70]
        (Cleft) circle (\radiusOuter); % left outer
    \draw[ultra thick, red!70]
        (Cright) circle (\radiusInner); % right inner
    \draw[ultra thick, red!70]
        (Cright) circle (\radiusOuter); % right outer
    \draw[ultra thick, red!70] (7.625,10) -- (7.625,-10);
    \draw[ultra thick, red!70] (10.625,10) -- (10.625,-10);

  \end{tikzpicture}

\end{column}
\end{columns}


\end{beamercolorbox}
\end{column}

%% ─ TRACKING ──────────────────────────────────────────────────────────────────
\begin{column}{0.50\textwidth}
  \begin{beamercolorbox}[rounded=true,sep=0.30cm,wd=\linewidth]{block title}
    \large\bfseries Tracking Controller
  \end{beamercolorbox}

  \begin{beamercolorbox}[rounded=true,sep=0.38cm,wd=\linewidth]{block body}

    \vspace{0.3cm}
    \centering
    %% ── MPC Description ────────────────────────────────────────────────────
    \begin{tcolorbox}[pbox]
    {\color{primary}\bfseries Motion Planning Overview}\\[3pt]
    \small
    \textbf{Linearised bicycle model MPC}\\
    Discrete-time model based on Linearised dynamic bicycle model:
    $[e_y,\ \dot e_y,\ e_\psi,\ \dot\psi]$ with steering input $\delta$.

    \textbf{QP formulation (unconstrained)}\\
    Precomputed prediction matrices for multiple velocities allow solving
    $U^{\ast} = -H^{-1} g$
    via Cholesky factorisation at runtime for real-time performance.

    \textbf{Path-relative state space}\\
    Vehicle pose is projected onto a spline track to compute states and the desired state is projected to go along the spline at the current speed.\\[4pt]
    Desired state = $[e_y = 0,\ \dot e_y = 0, \ e_\psi = 0,\ \dot\psi = \kappa \cdot v]$\\
    \textbf{GG-ellipse speed planner}\\
    Curvature-based speed limits are refined using a forward/backward pass
    over a friction ellipse to enforce combined lateral/longitudinal limits.

    \textbf{Longitudinal control}\\
    PID controler based on the desired speed of the GG-Plot and the current speed from odometry. The GGplot results is refined by using lookahead based on the current speed to slow down before corners.
    \end{tcolorbox}


    \vspace{0.3cm}
    \raggedright
    %% ── Diagrams side by side ───────────────────────────────────────────────
    \noindent
    \hspace{0.01\linewidth}
    \begin{minipage}[t]{0.35\linewidth}
      \centering
      %% GG-Plot
      \begin{tikzpicture}[scale=3.45]
        % Axes
        \draw[-{Stealth},accentgray,line width=0.9pt](-1.9,0)--(2.1,0)
          node[right,font=\small,accentgray]{$a_y$};
        \draw[-{Stealth},accentgray,line width=0.9pt](0,-1.9)--(0,2.1)
          node[above,font=\small,accentgray]{$a_x$};
        % Axis ticks 
        \draw[accentgray,line width=2.3pt] ( 1.71,0.06) -- ( 1.71,-0.06);
        \draw[accentgray,line width=2.3pt] (-1.71,0.06) -- (-1.71,-0.06);
        \draw[accentgray,line width=2.3pt] (0.06, 1.21) -- (-0.06, 1.21);
        \draw[accentgray,line width=2.3pt] (0.06,-1.21) -- (-0.06,-1.21);
        \node[accentgray,font=\scriptsize] at ( 1.9,-0.18){$a_{y,\max}$};
        \node[accentgray,font=\scriptsize,anchor=east] at (0.58, 1.25){$a_{x,\max}$};
        \draw[primary,dashed,line width=1pt] (0.73,1.08) -- (0.73,0);
        \draw[primary,dashed,line width=1pt] (0.73,1.08) -- (0,1.08);
        % GG friction circle
        \draw[secondary,line width=2.2pt,dashed] (0,0) ellipse (1.7 and 1.2);
        \fill[secondary!10,opacity=0.5] (0,0) ellipse (1.7 and 1.2);
        % Shade interior
        % Several speed-limit operating points at different corner radii
        % Highlight main op point
        \fill[primary](0.73,1.08) circle(0.04);
        \draw[-{Stealth},primary,line width=1.4pt](0,0)--(0.73,1.08);
        \node[primary,font=\scriptsize\bfseries,anchor=south west] at (0.92,1.22)
          {op.\ pt};
        % v_max annotation — right side
        %\draw[<->,red!60!black,line width=1.0pt](1.65,0.0)--(1.65,-1.50);
        %\node[red!60!black,font=\scriptsize\bfseries,anchor=west] at (1.68,-0.75)
        %  {$v_{\max}$};
        % Caption
        \node[accentgray,font=\small\bfseries] at (0.1,-2.2)
          {GG-plot $\to$ $v_{\max}(s)$};
      \end{tikzpicture}
    \end{minipage}%
    \begin{minipage}[t]{0.49\linewidth}
      \centering
      %% Lateral MPC track diagram
      \begin{tikzpicture}[scale=4.55]
        % Track boundaries (cones)
        \foreach \x/\y in {-2.8/0.95,-1.8/1.15,-0.8/1.25,0.2/1.15,1.2/1.05,2.2/0.85}
          {\fill[primary](\x,\y) circle(0.10);}
        \foreach \x/\y in {-2.8/-0.75,-1.8/-0.55,-0.8/-0.55,0.2/-0.55,1.2/-0.65,2.2/-0.75}
          {\fill[secondary](\x,\y) circle(0.10);}
        % Reference path
        \draw[secondary,line width=2.2pt,dashed]
          (-3,0) to[out=5,in=175] (-0.5,0.42) to[out=-5,in=185] (2.5,-0.1);
        % Vehicle body (rect)
        \fill[darkprimary,rounded corners=2pt, rotate=20]
          (-2.92,0.65) rectangle (-2.18,1.13);
        % Horizon preview dots
        \foreach \x/\y [count=\i] in {-2.0/0.21,-1.5/0.34,-1.0/0.42,-0.5/0.42,0.0/0.35} {
          \coordinate (P\i) at (\x,\y);
          \fill[primary!65] (P\i) circle(0.03);
        }
        \draw[->, ultra thick, red!65] (-2.36,0.08) -- (P1);
        \foreach \i in {1,...,4} { % 4 because we have 5 points
            \draw[->, ultra thick, red!65] (P\i) -- (P\the\numexpr\i+1\relax);
        }
        % MPC predicted path
        % Horizon brace
        \draw[<->,primary!80,line width=0.9pt](-2.5,0.75)--(0.0,0.75);
          \node[primary,font=\small\bfseries] at (-1.25,0.90){\textbf{Horizon} $N$};
        % Lateral error
        %\draw[<->,red!60!black,line width=1.0pt](-2.5,-0.50)--(-2.5,0.05);
        %\node[red!60!black,font=\scriptsize\bfseries,anchor=west] at (-2.38,-0.22)
        %  {$e_y$};
        % Heading error arc
        %\draw[-{Stealth},orange,line width=1.1pt]
        %  (-2.5,0.10) arc[start angle=180,end angle=153,radius=0.65];
        %\node[orange,font=\scriptsize\bfseries,anchor=east] at (-3.05,0.55)
        %  {$e_\psi$};
        % Caption
        \node[accentgray,font=\small\bfseries] at (0,-1.25)
          {Lateral MPC: predictive horizon};
      \end{tikzpicture}
    \end{minipage}

    \vspace{0.3cm}

    %% ── Rationale ───────────────────────────────────────────────────────────
    \begin{tcolorbox}[ybox]
      \small\color{darkprimary}
      \textbf{Rationale:} Decoupling lateral and longitudinal control allows us to solve two simpler problem
      rather than one complex one. A linear MPC gives a reliable and robust longitudonal control while
      being much simpler than a nonlinear one. The GG-ellipse speed planner allows us to enforce combined 
      slip limits without needing a complex tire model.
    \end{tcolorbox}

  \end{beamercolorbox}

\end{column}
%% ─ SIMULATOR ─────────────────────────────────────────────────────────────────
\begin{comment}
\begin{column}{0.355\textwidth}
\begin{beamercolorbox}[rounded=true,sep=0.30cm,wd=\linewidth]{block title}
  \large\bfseries Simulation Environment
\end{beamercolorbox}
\begin{beamercolorbox}[rounded=true,sep=0.38cm,wd=\linewidth]{block body}

\begin{columns}[t]
\begin{column}{0.48\linewidth}
  \begin{tcolorbox}[pbox]
    {\color{primary}\bfseries Project Chrono Vehicle Model}\\[3pt]
    \small
    High-fidelity physics simulation:\\
    $\bullet$ \textbf{Suspension} model\\
    $\bullet$ \textbf{Steering} kinematics\\
    $\bullet$ \textbf{Drivetrain} model\\
    $\bullet$ \textbf{Wheels \& tyres:}\\
    \quad load distribution,\\
    \quad longitudinal \& lateral model
  \end{tcolorbox}

  \vspace{0.3cm}
  \begin{tcolorbox}[pbox]
    {\color{primary}\bfseries Custom LiDAR Simulator}\\[3pt]
    \small
    Built in-house to match\\
    Ouster OS1-32 characteristics:\\
    $\bullet$ 32-channel ray casting\\
    $\bullet$ Gaussian range noise\\
    $\bullet$ Beam dropout model\\
    $\bullet$ Correct FoV \& scan rate
  \end{tcolorbox}
\end{column}
\begin{column}{0.48\linewidth}
  % Chrono sim diagram
  \centering
  \begin{tikzpicture}[node distance=0.52cm,scale=0.88]
    \tikzstyle{cn}=[fill=primary,text=white,rounded corners=3pt,
      minimum width=3.8cm,minimum height=0.55cm,align=center,font=\small\bfseries]
    \tikzstyle{cy}=[fill=secondary,text=darkprimary,rounded corners=3pt,
      minimum width=3.8cm,minimum height=0.55cm,align=center,font=\small\bfseries]
    \node[cy](ph){Project Chrono Physics};
    \node[cn,below=0.38cm of ph](lidar){Custom LiDAR Sim};
    \node[cn,below=0.38cm of lidar](imu){IMU Emulator};
    \node[cy,below=0.38cm of imu](ros){ROS2 Interface};
    \node[cn,below=0.38cm of ros](stack){Autonomous Stack};
    \node[cy,below=0.38cm of stack](cmd){Actuator Cmds};
    \foreach \a/\b in {ph/lidar,lidar/imu,imu/ros,ros/stack,stack/cmd}
      {\draw[-{Stealth[length=5pt]},primary,line width=1.2pt](\a)--(\b);}
    \draw[darkprimary,dashed,line width=0.8pt]
      (cmd.east)--++(0.9,0)|-(ph.east);
    \node[accentgray,font=\scriptsize,rotate=90]at(2.65,-2.7){feedback};
  \end{tikzpicture}

  \vspace{0.3cm}
  \begin{tcolorbox}[ybox]
    \small\color{darkprimary}
    \textbf{Benefit:} Full stack tested in simulation before any physical trial, with realistic tyre and sensor behavior
  \end{tcolorbox}
\end{column}
\end{columns}

\vspace{0.3cm}
% Sim track diagram
\centering
\begin{tikzpicture}[scale=0.82]
  \draw[primary!20,line width=8pt,rounded corners=22pt]
    (-4.5,0)--(-3.2,1.8)--(-1.0,2.8)--(1.5,2.8)--(3.2,1.8)--(4.0,0)
    --(3.2,-1.8)--(1.5,-2.5)--(-1.0,-2.5)--(-3.2,-1.8)--cycle;
  \draw[secondary!35,line width=3.5pt,rounded corners=14pt]
    (-2.9,0)--(-2.1,1.3)--(-0.7,2.0)--(1.2,2.0)--(2.3,1.3)--(2.7,0)
    --(2.3,-1.3)--(1.2,-1.8)--(-0.7,-1.8)--(-2.1,-1.3)--cycle;
  % Cones
  \foreach \x/\y in {-4.5/0,-3.2/1.8,-1.0/2.8,1.5/2.8,3.2/1.8,4.0/0,
                      3.2/-1.8,1.5/-2.5,-1.0/-2.5,-3.2/-1.8}
    {\fill[primary](\x,\y)circle(0.13);}
  \foreach \x/\y in {-2.9/0,-2.1/1.3,-0.7/2.0,1.2/2.0,2.3/1.3,2.7/0,
                      2.3/-1.3,1.2/-1.8,-0.7/-1.8,-2.1/-1.3}
    {\fill[secondary](\x,\y)circle(0.11);}
  % Car
  \fill[darkprimary](-0.2,0.1)circle(0.18);
  % Centerline
  \draw[darkprimary,line width=1.4pt,dashed,rounded corners=12pt]
    (-1.8,0)--(-1.1,1.0)--(0,1.4)--(1.1,1.0)--(1.7,0)
    --(1.1,-1.0)--(0,-1.4)--(-1.1,-1.0)--(-1.8,0);
  % LiDAR rays from car
  \foreach \a in {-40,-20,0,20,40}
    {\draw[secondary,opacity=0.5,line width=0.8pt](-0.2,0.1)--+(\a:2.5);}
  \node[accentgray,font=\small]at(0,-3.1){Simulation track with LiDAR rays};
\end{tikzpicture}

\end{beamercolorbox}
\end{column}
\end{comment}

\end{columns}
\vspace{0.3cm}
\end{frame}
\end{document}
